\section{Conclusiones}
\label{sec:conclusiones}

El protocolo cumple las garantías que el enunciado pide, y lo sostienen dos
clases independientes de evidencia: un verificador externo a la
implementación, que no halló violaciones de conservación, unicidad,
integridad ni FIFO en ninguna corrida, y una suite con intercalados
difíciles de provocar a mano ---cierre en plena actividad, incorporaciones
durante el cierre, cuatro finalizadores compitiendo--- bajo sanitizers sin
reportes. Las dos son necesarias: los sanitizers no buscan carreras y el
verificador no detecta errores de memoria.

Del experimento, el resultado central es que subir la capacidad del buffer
no es una mejora sino un intercambio: aumenta el caudal solo cuando ninguno
de los dos lados es el cuello de botella, mientras la latencia sube siempre
y, en régimen saturado, deja de medir el protocolo para volverse la
capacidad dividida por la tasa de drenaje.

Y queda nítida la frontera entre lo que garantiza el algoritmo ---exclusión,
conservación, orden y terminación--- y lo que concede el sistema operativo,
que decide quién de los que esperan avanza y cuándo. Ninguna propiedad de
corrección depende de esa frontera, pero varios números sí. El alcance
tampoco se amplía al final: no hay recuperación ante la caída arbitraria de
un proceso, no existe herramienta que busque carreras entre procesos
separados sobre memoria compartida, y las mediciones vienen de un binario
sin optimizar sobre un hipervisor.
